\newpage
\section{Characterized IP Pools Used in the Evaluation}
\label{app:ip_pools}

This appendix reproduces the three IP pools used by the Level~2
exploration in the experiments of Chapter~\ref{sec:results}: the INT8
pool (\path{ip_pool_characterized_65nm.yaml}), the INT16 pool
(\path{ip_pool_int16_65nm.yaml}) and the FP16 pool
(\path{ip_pool_fp16_65nm.yaml}). The INT16 pool is used for the two
AlexNet experiments, the INT8 pool for the MobileNet/Eyeriss~v2
comparison and the FP16 pool for the VGG-16/SAURIA comparison. The
processing elements and memories they contain are summarized in
Tables~\ref{tab:characterized_pes}
and~\ref{tab:characterized_memories}; the meaning of every field is
described in Section~\ref{subsec:ip_pool} and
Appendix~\ref{app:talos_requirements}.

The files are included exactly as consumed by Talos. All IPs were
characterized with Cadence Genus after mapping to the TSMC 65~nm
standard-cell library at 200~MHz, 1.2~V, 25~$^\circ$C and the NCCOM
corner (Section~\ref{subsec:ip_pool}). Areas are expressed in mm$^2$,
power in W and delays in ns. The three pools share the same twelve
register files and three global buffers; only the compute entries
change. The \texttt{dram\_platform\_512b} entry of the INT8 and INT16
pools is the shared synthetic platform proxy described in
Appendix~\ref{app:talos_requirements}; the FP16 pool omits it and lets
the full flow inject the same proxy. The pools are stored in the
thesis repository under \path{ip_pools/} and correspond to the
\path{configs/} directory of the Talos repository.

\lstdefinestyle{yamlpool}{
    language={},
    basicstyle=\scriptsize\ttfamily,
    columns=fixed,
    basewidth=0.5em,
    keepspaces=true,
    breaklines=true,
    breakindent=0pt,
    showstringspaces=false,
    frame=single,
    rulecolor=\color{black!40},
    backgroundcolor=\color{background},
    numbers=none,
    captionpos=t,
    nolol=true,
    upquote=true,
    sensitive=true,
    morecomment=[l]{\#},
    commentstyle=\color{black!60}\itshape,
    morestring=[b]",
    stringstyle=\color{black},
    literate={\$}{{\$}}1
}

\subsection{Annotated Example of an IP Entry}
\label{app:ip_pool_example}

Listing~\ref{lst:ip_pool_annotated} shows one processing element and
one register file of the INT16 pool with a comment on every field.
The same structure is repeated for every entry of the pools
reproduced in the following sections. Fields under
\texttt{metadata} are not interpreted by the generic IP model; Level~2
reads from them only the values it needs (\texttt{macs\_per\_cycle},
\texttt{accesses\_per\_cycle}, and the numerical format of the PE),
while the remaining values are kept as documentation of the
characterization.

\begin{lstlisting}[style=yamlpool,nolol=false,label={lst:ip_pool_annotated},caption={Annotated INT16 pool entries: one PE and one register file.}]
technology_nm: 65                # technology node shared by the pool;
                                 # also used by the Level 1 CACTI calibration
ips:
  - id: pe_sauria_int16          # unique identifier of the IP
    type: pe                     # pe | register_file | global_buffer | dram
    area: 0.00509364             # area of one instance [mm2]
    throughput: 1.0              # operations per cycle of the generic model
    delay: 4.803                 # critical path [ns]
    fmax_mhz: 208.20320632937748 # maximum frequency derived from the delay
    metadata:
      precision_bits: 16         # compared with the ONNX workload operands
      activation_bits: 16
      weight_bits: 16
      accumulator_bits: 48
      output_bits: 48
      signed: true
      data_type: integer
      numeric_format: int16
      macs_per_cycle: 1          # used to check the Level 1 mapping
      pipeline_latency_cycles: 3
      dynamic_energy_per_mac_pj: 5.281205   # (active - idle) / f_ref
    power_model:                 # idle/active characterization
      source: Cadence Genus mapped standard-cell characterization
      activity_method: vectorless_explicit
      reference_frequency_mhz: 200.0        # operating point of the
      p_idle_w: 0.000593089                 # characterization; all IPs
      p_active_w: 0.00164933                # of one implementation must
      voltage_v: 1.2                        # share it
      temperature_c: 25.0
      corner: NCCOM

  - id: mem_bsg_1rw_8x64
    type: register_file
    area: 0.00531648
    throughput: 1.0
    delay: 1.143
    fmax_mhz: 874.8906386701664
    capacity_bits: 512           # must cover the abstract memory capacity
    bandwidth_bits: 64           # must cover the abstract memory bandwidth
    metadata:
      accesses_per_cycle: 1      # compared with the ZigZag memory traffic
      ports: 1RW
      read_energy_pj: 0.7066
      write_energy_pj: 3.6565
    power_model:                 # same operating point as the PE
      source: Cadence Genus mapped standard-cell characterization
      activity_method: vectorless_explicit
      reference_frequency_mhz: 200.0
      p_idle_w: 0.00111907
      p_active_w: 0.00174303
      voltage_v: 1.2
      temperature_c: 25.0
      corner: NCCOM
\end{lstlisting}

In the INT16 and FP16 files the operating-point fields are written
once with a YAML anchor (\texttt{\&characterization}) and reused by
the remaining entries with a merge key (\texttt{<{}<: *characterization}); only the idle and active powers are overridden.
The INT8 file lists every field explicitly. Both forms are equivalent
after parsing.

\clearpage
\subsection{INT8 Pool}
\label{app:ip_pool_int8}

\lstinputlisting[style=yamlpool,caption={\path{ip_pool_characterized_65nm.yaml}: INT8 processing elements, shared memories and platform DRAM proxy.}]{ip_pools/ip_pool_characterized_65nm.yaml}

\clearpage
\subsection{INT16 Pool}
\label{app:ip_pool_int16}

\lstinputlisting[style=yamlpool,caption={\path{ip_pool_int16_65nm.yaml}: INT16 processing elements, shared memories and platform DRAM proxy.}]{ip_pools/ip_pool_int16_65nm.yaml}

\clearpage
\subsection{FP16 Pool}
\label{app:ip_pool_fp16}

\lstinputlisting[style=yamlpool,caption={\path{ip_pool_fp16_65nm.yaml}: FP16 processing elements and shared memories.}]{ip_pools/ip_pool_fp16_65nm.yaml}
